\section{Introduction}
Generating a group of responses defines a distribution over comparisons as well as
a distribution over answers. Two samplers can preserve the marginal policy and
have the same complete distribution of binary success counts while producing
oppositely directed expected policy updates. Coverage therefore does not determine
what a group-relative estimator teaches the policy.

Dependence-induced baseline bias is established prior work. QuasiMoTTo studies
marginal-preserving sequence sampling, its RLOO bias, and pairwise correction
\citep{li2026quasimotto}; CARMS and ARMS develop coupled score-function estimators
\citep{dimitriev2021carms,dimitriev2021arms}. Our target is the interaction between
sampling dependence and the policy's learnable reward directions, together with
a predictive causal test that preserves the reward-count explanation.

The analysis uses classical reversible-kernel geometry and score-space projections
\citep{sherlock2018reversible,sutton1999policy,kakade2001natural}. The mathematical
ingredients are not claimed as new. Their application identifies a precise boundary:
dependence is a reward-independent positive natural preconditioner exactly when
the score subspace is invariant under the partner operator. Otherwise, a reward
component invisible to the original policy gradient can enter the coupled update.

For binary rewards, we further separate count-associated scale from the identities
that receive positive or negative advantages. Permuting successful identities over
the same successful slots, and failed identities over failed slots, preserves every
group's count and all stored trajectories. Its analytic average removes additional
permutation noise. The proposed empirical contribution is to predict the resulting
functional and verifier-reward changes from independent data. No such pretrained-model
claim is made in this initial version.

\paragraph{Scope of the draft.}
We provide the compatibility characterization, sampler consequences, the binary
count/identity decomposition, and finite-policy checks. The planned empirical
contribution is count-preserving causal identification supported by independent
IID forecasts and held-out reward consequences. Cross-group baselines serve as
an identification instrument, following established baseline principles
\citep{greensmith2004variance,chung2020beyond}.
